\section{Conclusion}
\label{sec:conclusion}

We evaluated four fast-training techniques from a state-of-the-art CNN recipe on a Vision Transformer trained on \cifar{}, under both a weak recipe ($61.7\%$ baseline) and a strong recipe ($81.4\%$ baseline).

Under the weak recipe, representation-quality tricks transfer positively (whitening $+1.3$\pp{}, alternating flip $+0.8$\pp{}) while architecture-specific tricks are harmful (Dirac $-7.9$\pp{}, ScaleBias+Lookahead $-16.0$\pp{}).
Under the strong recipe the picture partially reverses, with whitening flipping to $-4.4$\pp{} and soft Dirac becoming the only net-positive trick ($+0.4$\pp{}).

The central lesson is that trick transferability is a \emph{trick~$\times$~recipe} interaction, not a fixed architectural property.
Frozen whitening is beneficial only when training augmentation preserves the patch statistics of the whitening basis; Mixup and CutMix destroy that match.
Dirac init is harmful when augmentation is too weak to quickly break attention symmetry, and mildly helpful when it is strong enough.
Lookahead conflicts with any monotonically decaying LR schedule regardless of architecture.

A controlled CNN masked-reconstruction experiment directly confirmed the whitening mechanism: trainable and frozen whitening perform identically when inputs are unaugmented, isolating the input-distribution mismatch as the sole cause of the reversal.

\paragraph{Future work.}
The most direct next steps are: (1) learnable (unfrozen) whitening, which should recover whitening's benefit under strong augmentation; (2) decayed soft Dirac ($\alpha$ annealed to 0), which may turn further positive; (3) Lookahead paired with a flat LR to remove the schedule conflict; and (4) an all-methods combination using all three fixes, to test whether the corrected tricks compose additively.
